\section*{Part VII: The Continuing Case: Bayes Institute as a Capability-Building Institution}

\subsection*{91. The Institutional Mission Is Not the 1st Course}

Bayes Institute's enduring philosophy is that \textbf{humanity's greatest promise is what people can become}. Its mission is to \textbf{create the conditions in which human possibility becomes human capability}. The institutional campaign is \textbf{BECOME MORE.}

Those commitments express an aspiration and a standard of responsibility. They do not decide the 1st product, buyer, duration, technology stack, or price. A mission about capability becomes commercially credible only when an institution can demonstrate a bounded capability it reliably helps actual learners develop.

The operating problem is therefore narrower than the philosophy: \textbf{Can an extremely lean, founder-led STEM institute help a specifically prepared adult make a substantial, independently demonstrable improvement through a learning experience that learners value and the institute can repeat?} This remains a hypothesis, not a confirmed Bayes result. The answer must come from experiments and real learner evidence.

\subsection*{92. 3 Potential Customers Imply 3 Different Transformations}

Bayes has considered 3 starting business routes. Their instructional mechanisms overlap, but their unit of value differs.

\begin{table*}[t]
\centering
\small
\caption{3 Potential Customers and Their Distinct Transformations}
\begin{tabularx}{\textwidth}{p{0.20\textwidth}p{0.28\textwidth}p{0.26\textwidth}X}
\toprule
\textbf{Possible Customer} & \textbf{Valuable Transformation} & \textbf{What Must Be Observed} & \textbf{What Does Not Follow Automatically} \\
\midrule
Working technical adult & Greater practical capability on defined engineering tasks & Independent work, feedback response, delayed and varied tasks & Promotion, placement, or credential prestige \\
College or university & Better student learning integrated into an academic setting & Student baseline, delivery quality, faculty integration, assessment & Signed multi-year partnership or improved placements \\
Enterprise team & Improved workflow-relevant technical behavior & Team performance on approved tasks and actual application & Attributable financial ROI from course attendance alone \\
\bottomrule
\end{tabularx}
\end{table*}

A program built around 1 transformation may not satisfy the other 2 without redesign. Selling all 3 as ``AI upskilling'' obscures critical differences in stakes, prerequisites, purchasing processes, and what counts as success.

\subsection*{93. A 1st Learner Hypothesis, Not a Presumed Product-Market Fit}

Consider a hypothetical target: working software engineers who can write production code but lack confidence and tested capability in evaluating and shipping AI features. They are not simply ``people interested in AI.'' They have a relevant technical base and a consequential task family.

The proposed educational gap is not necessarily lack of access to introductory lectures. It may be lack of \textbf{expert critique, diagnostic practice, and structured opportunities to defend decisions under production-style constraints}. This is plausible given the founder's expertise, but customer discovery must establish whether the problem is acute enough to earn money and protected time.

A small diagnostic sprint might test whether learners value that kind of feedback. It would \textbf{not} validate a 6-month program, a job-market outcome, or long-term institutional credibility. Those require stronger and different evidence.

\subsection*{94. Define the Learner's Starting State Before Naming the Program}

The entry state for the hypothetical engineer should be assessed across distinct dimensions: ability to read and modify unfamiliar code; familiarity with data and model evaluation; reliability and testing practices; reasoning about trade-offs; and ability to explain uncertainty.

A 20-minute self-report form could miss consequential weaknesses. Better evidence might include a short unfamiliar code-reading task, a flawed evaluation design, and a written decision about whether to ship a small AI component. Different students can be strong and weak in different dimensions.

Bayes should define \textbf{minimum admissible prerequisites} and \textbf{bridgeable gaps}. Admitting every interested person may increase apparent demand while making a rigorous shared learning outcome expensive or impossible. Conversely, admitting only learners who already possess the target competence makes later success easy to market but weak as evidence of transformation.

\subsection*{95. A Capability Claim That Can Actually Be Tested}

Suppose the target is: \emph{A technically prepared learner can evaluate an unfamiliar AI application, diagnose important failure modes, propose and test improvements, and defend the remaining risk under explicit time and tool constraints.}

A more specific formulation might require the learner to handle a new system, construct or critique a representative test set, compare a simple baseline with an improved approach, inspect errors, explain cost and reliability trade-offs, and respond to a surprise change request.

This is still too broad for a 4-week experiment. The founder should select 1 consequential sub-capability, perhaps \textbf{detecting and repairing misleading AI evaluation practices}, and design a narrow learning loop around it. A smaller, complete transformation is more informative than a large collection of introductory topics.

\subsection*{96. Build the 1st Experience Around Error Mechanisms}

For a diagnostic-and-improvement sprint, select authentic but privacy-safe failure cases. 1 could involve evaluation-data leakage; another could involve retrieval failure obscured by flattering aggregate metrics; a 3rd could involve a misleading baseline caused by input selection.

Begin with an unseen diagnostic and record reasoning, not merely accuracy. Then use concise explanations and worked examples to expose the key mechanisms. Require learners to attempt varied cases, receive targeted feedback, and revise. Finish with a new failure case requiring an independent diagnosis and defense.

This is \textbf{a hypothetical instructional design}, not an established Bayes offer. Its purpose is to demonstrate whether the founder can create a measurable, valuable learning experience without 1st constructing a large platform or recording months of videos.
